% Data appendix: ACS same-age and same-sex sibling designs for housing.
% Self-contained component. Requires amsmath, booktabs, graphicx, natbib
% (all loaded by the main draft). The figure directory can be overridden by
% defining \acsivfigdir before this file is input.
\providecommand{\acsivfigdir}{../../output/acs_fertility_iv/data_appendix/figures}

\section{Housing and Additional Children: Evidence from Sibling Composition in the ACS}
\label{app:acs_iv}

This appendix asks how the housing a young mother occupies changes when she has
an additional child, using two features of her children that are plausibly
unrelated to her housing preferences except through family size. The first is
whether her two oldest children are recorded at the same age, which in most cases
reflects a twin birth. The second is whether her two oldest children have the same
sex, which \citet{angrist1998children} show raises the probability of a third
birth. Both sources of variation have a long history in the study of fertility and
labor supply. Housing is a different outcome, and for the sibling-sex design an
awkward one: children of the same sex can share a bedroom, so the composition of a
family can change the housing it needs without changing its size.

The two designs study different families at different margins, and they give
answers of opposite sign. Mothers whose two oldest children
are the same age live in about a third of a room more than comparable mothers, and
are about three percentage points more likely to own their home. Mothers whose two
oldest children have the same sex live in slightly \emph{fewer} rooms, a difference
almost entirely accounted for by bedrooms. Taken at face value, the corresponding
instrumental-variable estimate says that a third child reduces a family's dwelling
by more than a full room. We read this estimate as evidence that the exclusion
restriction fails for housing outcomes, not as an effect of children, and we
explain below why a small direct effect of sibling sex on housing is enough to
produce it.

\subsection{Data and the mother sample}
\label{app:acs_iv_data}

The data are the one-year samples of the American Community Survey (ACS) for
2005--2019, covering the fifty states and the District of Columbia, as distributed
by IPUMS USA. Each survey year is a separate cross-section of households. A
person is identified by the survey year, the sample, the household serial number,
and the person's number within the household; a household is identified by the
first three of these. No person is followed across years, so every quantity below
is measured once, at the interview.

The ACS does not ask women how many children they have borne; it asks only whether
a woman gave birth in the past twelve months, and it records no birth histories. We therefore reconstruct each
woman's children from the household roster. IPUMS assigns to each person the
person number of a probable mother living in the same household, using
relationship, age, and marital information \citep{ipumsusa_vardocs}. These links
identify social as well as biological relationships: a stepmother or adoptive
mother is linked in the same way as a birth mother. We keep every link in which the
child's age is nonnegative and below the mother's age. The household relationship
variable, which records each person's relationship to the householder rather than
to the mother, cannot separate biological from other links and is not used to
filter them. Among all valid links in the 2005--2019 records, before any
restriction on ages, 91 percent of linked children are recorded as the
householder's child and 4 percent as the householder's grandchild, the latter
typically a young mother living in her own parents' home. Only 1,692 links fail
because the recorded child is not younger than the recorded mother.

For each woman we count her linked coresident children, $n_i$, and order them by
age, $a_{i1}\ge a_{i2}\ge\cdots$, where ages are whole years at the interview. The
count $n_i$ is the number of her children living with her, not the number she has
borne. Children who have left home, died, or live with another parent are
missing, and a woman can be linked to a child she did not bear. Following
\citet[pp.~452--453]{angrist1998children}, we limit these errors by studying young
mothers: the base sample consists of women aged 21 to 35 at the interview whose
oldest linked child, if any, is younger than 18. Angrist and Evans additionally
drop mothers whose count of coresident children differs from their reported
children ever born and mothers whose children's ages or sexes were imputed by the
Census Bureau. The ACS does not collect children ever born, and our extract does
not carry imputation flags, so neither restriction is applied here.

Table~\ref{tab:acsiv_sample} reports the resulting samples. The 59,046,776 person
records in the source files fall to 46,373,936 once the sample is limited to the
2005--2019 one-year files. The base sample contains 4,024,238 women living in
3,783,862 households. It includes women with no linked child, who enter neither
design, and it counts women rather than mothers.

\begin{table}[htbp]
\centering
\footnotesize
\caption{Sample construction}
\label{tab:acsiv_sample}
\begin{tabular}{@{}lr@{}}
\toprule
 & Count \\
\midrule
\multicolumn{2}{@{}l}{\textit{Source records}} \\
ACS person records, fifty states and the District of Columbia & 59,046,776 \\
\quad in the 2005--2019 one-year samples & 46,373,936 \\
Women aged 21--35 with positive person weight and no linked child aged 18 or older & 4,024,238 \\
\quad households containing at least one such woman & 3,783,862 \\
\addlinespace
\multicolumn{2}{@{}l}{\textit{Same-age design}} \\
Mothers whose oldest linked child is aged 0--5 & 857,607 \\
\quad with two or more linked children & 325,171 \\
\quad whose two oldest linked children are the same age & 16,496 \\
\quad whose three oldest linked children are the same age & 388 \\
\addlinespace
\multicolumn{2}{@{}l}{\textit{Same-sex design}} \\
Women with two or more linked children & 1,132,295 \\
\quad whose two oldest linked children are the same age & 32,186 \\
Mothers meeting all criteria (text) & 656,986 \\
\quad two oldest children of the same sex & 329,974 \\
\qquad two boys & 173,813 \\
\qquad two girls & 156,161 \\
\quad three or more linked children & 148,775 \\
\bottomrule
\end{tabular}
\par\medskip
\begin{minipage}{0.95\linewidth}
\footnotesize\textit{Notes:} Counts of person records, unweighted. Nested entries are subsets of the entry above
them. In the same-sex design, the count of women whose two oldest children share an
age is taken over all women with two or more linked children, whatever the age of
the second child; the criteria for the estimation sample are listed in
Section~\ref{app:acs_iv_instruments}.
\end{minipage}
\end{table}

\subsection{Housing outcomes}
\label{app:acs_iv_outcomes}

We measure three features of the dwelling, all reported once per household and
assigned to every mother in it. The universe of each is housing units outside group
quarters.

\emph{Rooms} counts whole rooms used for living purposes, including kitchens and
excluding bathrooms, porches, balconies, foyers, halls, and half-rooms. IPUMS codes
valid counts from 1 to 27 and a top category of 30. In the 2009 file, 107 person
records carry the value 28, which the codebook does not define; they are treated as
missing. We cap the count at nine rooms because the Census Bureau imposed
a national top code of nine rooms through the 2007 survey
\citep{ipumsusa_vardocs}. Without the cap, a family's measured size would depend on
the survey year in which it was interviewed.

\emph{Bedrooms} counts rooms that would be listed as bedrooms if the unit were for
sale or rent, whatever their current use. The IPUMS code minus one is the number of
bedrooms; valid codes run from 1 (no bedroom) to 22. We cap the count at five,
again because a national top code of five bedrooms applied through 2007. Bedrooms
are a component of rooms, and we report them to see which part of the dwelling
moves.

\emph{Ownership} equals one if the unit is owned or being bought, including units
held with a mortgage, and zero if it is rented, including units occupied without
payment of cash rent. Other codes are missing.

Every mother in both estimation samples has a valid tenure code and a valid bedroom
count. One mother in the same-sex sample lacks a valid room count, which is why the
room samples in Table~\ref{tab:acsiv_pooled} are one observation smaller.

\subsection{The two instruments}
\label{app:acs_iv_instruments}

\paragraph{Same-age oldest pair.}
This design studies the move from one to two children among mothers whose first
child is young. The sample consists of mothers with at least one linked child whose
oldest linked child is aged 0 to 5. The instrument and the treatment are
\begin{equation}
Z_i^{A}=\mathbf{1}\{n_i\ge 2,\ a_{i1}=a_{i2}\},
\qquad
D_i^{A}=\mathbf{1}\{n_i\ge 2\},
\label{eq:acsiv_twin_def}
\end{equation}
so the instrument indicates that the two oldest children are recorded at the same
age and the treatment indicates a second coresident child. Every mother with
$Z_i^{A}=1$ has $D_i^{A}=1$ by construction. The event age is the oldest child's
age, $e_i=a_{i1}$, which measures the time since the first birth in whole years.

The instrument is a proxy for a twin first birth, not a twin indicator. Ages are
recorded in whole years and our extract carries no quarter of birth, so two
children born within a year of each other can share an age at the interview. A
step-sibling or adopted child of the same age produces the same record.
\citet[p.~454]{angrist1998children} measure the resulting error in the 1980 Census:
defining twins by age in years alone yields a second-birth twin rate 35 percent
higher than defining them by age and quarter of birth, and for that reason they
restrict their twins analysis to 1980. In our sample 1.9 percent of eligible mothers
have a same-age pair. Genuine twin births are themselves not random:
\citet[p.~775]{angrist2010multiple} note that twinning varies with maternal age and
race and that twin births change child spacing and child health.

\paragraph{Same-sex oldest pair.}
This design studies the move from two to three children. The sample consists of
mothers with at least two linked children whose two oldest children have different
ages, whose second-oldest child is aged 0 to 5, and whose two oldest children both
have a recorded sex. The instrument and treatment are
\begin{equation}
Z_i^{B}=\mathbf{1}\{s_{i1}=s_{i2}\},
\qquad
D_i^{B}=\mathbf{1}\{n_i\ge 3\},
\label{eq:acsiv_samesex_def}
\end{equation}
where $s_{ik}$ is the sex of the $k$th oldest linked child. The event age is the
second child's age, $e_i=a_{i2}$, so the clock starts at the second birth, when the
instrument is realized. The requirement that the two oldest children differ in age
removes the mothers who drive the first design and makes the order of the first
two children well defined; the order of the second and third children is not
checked, so when they share an age the child labeled second is arbitrary. The sample
is not restricted to mothers with exactly two children, and it does not condition
on the treatment. Every mother who satisfies the age and ordering criteria has a
recorded sex for both children.

The instrument predicts a third coresident child, not a second, and its first stage
is positive: parents of two boys or two girls are more likely to have another
child. At event age 0 the second child was born within the past year, but a third
child may already be present, so no event age in this design precedes treatment.

\subsection{Estimating equations and inference}
\label{app:acs_iv_equations}

For each design, outcome $Y_i$, and window of event ages, we estimate a reduced form
and a first stage on the same mothers:
\begin{equation}
Y_i=\delta Z_i+W_i'\gamma+u_i,
\qquad
D_i=\pi Z_i+W_i'\kappa+v_i .
\label{eq:acsiv_rf_fs}
\end{equation}
The reduced-form coefficient $\delta$ is the difference in housing between mothers
with and without the sibling feature, and $\pi$ is the corresponding difference in
the probability of the additional child. The controls $W_i$ are complete sets of
indicators for the mother's age at the event (her age at the interview minus
$e_i$), her race, the survey year, and the event age. We do not control for the
mother's current marital status, education, or state of residence, all of which can
respond to fertility; \citet[p.~465, footnote~12]{angrist1998children} exclude
education and husband's earnings for the same reason. Regressions are weighted by
the mother's ACS person weight. Standard errors are clustered by household, because
several eligible women can live in one household, and each outcome is estimated on
the mothers for whom it is observed.

With one instrument and one endogenous variable, two-stage least squares is the
ratio of the two coefficients:
\begin{equation}
\hat\beta=\hat\delta/\hat\pi .
\label{eq:acsiv_ratio}
\end{equation}
If the instrument is as good as randomly assigned given $W_i$, affects housing only
through the additional child, and moves every mother's fertility in the same
direction, $\beta$ is the average effect of the additional child for mothers whose
fertility responds to the instrument. For the first design these are mothers who,
without a same-age pair, would have one coresident child at that event age. For the
second they are mothers who have a third child because their first two share a sex.
The two groups differ in age, in family size, and in the margin they cross, so the
two estimates answer different questions.

Our baseline does not include separate controls for the sex of the first and second
child. \citet[pp.~460--461]{angrist1998children} add them because the same-sex
indicator is an interaction of the two children's sexes: with slightly more boys
than girls born, it is mildly correlated with each child's sex, and a child's sex
may affect parents for reasons unrelated to family size. The same authors report
that their estimates are virtually unchanged by these controls (pp.~466--467); we do
not report that specification here.

We report three measures of uncertainty. Confidence intervals are the estimate plus
or minus 1.96 clustered standard errors. The first-stage $F$ statistic is
$(\hat\pi/\widehat{\mathrm{se}}(\hat\pi))^2$ from the clustered covariance; with a
single instrument and a single endogenous variable, it coincides with the
Kleibergen--Paap statistic \citep{andrews2019weak}. Because a large $F$
guards against weak-instrument bias but says nothing about the exclusion
restriction, we also invert the Anderson--Rubin test
\citep{anderson1949estimation}, which \citet{andrews2019weak} recommend
reporting whenever there is a single instrument. For each candidate value $\beta_0$
we regress $Y_i-\beta_0D_i$ on $Z_i$ and $W_i$ and test whether the coefficient on
$Z_i$ is zero at the 5 percent level with the same clustered covariance. The
candidates form a grid from $-3$ to $3$ in steps of $0.1$ for rooms and bedrooms,
and from $-50$ to $50$ percentage points in steps of $2$ for ownership. The tables
report the smallest and largest grid values that are not rejected. With one
instrument, the Anderson--Rubin set is the set of $\beta_0$ satisfying a quadratic
inequality, and it is a bounded interval exactly when the first-stage $F$ exceeds
the critical value of 3.84; every first-stage $F$ in our tables does. The exact
endpoints lie between the last accepted and the first rejected grid value. Where one grid value is accepted, as for rooms in
the pooled same-age design, the set lies within one grid step of that value; it is
not a point. For bedrooms in the same design no grid value is accepted: the set is
narrower than the grid step and lies between two adjacent grid values, in line with
the conventional interval of $[0.305,0.361]$.

\subsection{Results}
\label{app:acs_iv_results}

Table~\ref{tab:acsiv_pooled} reports the estimates pooled over event ages 0 to 5,
and Table~\ref{tab:acsiv_ages} reports them at event ages 3 and 5. Figure~\ref{fig:acsiv_rf} plots all eighteen reduced
forms. We report every design, outcome, and window that was estimated.

\begin{table}[htbp]
\centering
\footnotesize
\caption{Pooled estimates, event ages 0--5}
\label{tab:acsiv_pooled}
\setlength{\tabcolsep}{4pt}
\begin{tabular}{@{}lrrrrrrrl@{}}
\toprule
 & & & Instrument & Reduced & First stage & First-stage & & AR accepted \\
Outcome & Mothers & Households & equal to one & form & (pp) & $F$ & 2SLS & grid values \\
\midrule
\multicolumn{9}{@{}l}{\textit{A. Same-age oldest pair, event age 0--5}} \\
Rooms & 857,607 & 855,366 & 16,496 & 0.329 & 66.21 & 77,396 & 0.496 & \{0.5\} \\
 & & & & (0.017) & (0.24) & & (0.026) & \\
Bedrooms & 857,607 & 855,366 & 16,496 & 0.221 & 66.21 & 77,396 & 0.333 & none \\
 & & & & (0.010) & (0.24) & & (0.014) & \\
Owner (pp) & 857,607 & 855,366 & 16,496 & 3.11 & 66.21 & 77,396 & 4.69 & [4, 6] \\
 & & & & (0.46) & (0.24) & & (0.70) & \\
\addlinespace
\multicolumn{9}{@{}l}{\textit{B. Same-sex oldest pair, event age 0--5}} \\
Rooms & 656,985 & 656,159 & 329,974 & $-$0.038 & 3.27 & 715 & $-$1.149 & [$-$1.4, $-$0.9] \\
 & & & & (0.005) & (0.12) & & (0.164) & \\
Bedrooms & 656,986 & 656,160 & 329,974 & $-$0.034 & 3.27 & 715 & $-$1.042 & [$-$1.2, $-$0.9] \\
 & & & & (0.003) & (0.12) & & (0.096) & \\
Owner (pp) & 656,986 & 656,160 & 329,974 & $-$0.25 & 3.27 & 715 & $-$7.50 & [$-$16, 0] \\
 & & & & (0.15) & (0.12) & & (4.51) & \\
\bottomrule
\end{tabular}
\par\medskip
\begin{minipage}{0.97\linewidth}
\footnotesize\textit{Notes:} Each row reports the reduced form and first stage of
equation~\eqref{eq:acsiv_rf_fs} and the two-stage least-squares estimate of
equation~\eqref{eq:acsiv_ratio}, with clustered standard errors in parentheses.
The treatment is a second coresident child in panel~A and a third in panel~B. Rooms
and bedrooms are counts, capped at nine and five; ownership and the first stage are
in percentage points. ``AR accepted grid values'' gives the smallest and largest
values of the effect not rejected by the Anderson--Rubin test at the 5 percent
level on the grid described in the text; braces denote a single accepted value and
``none'' an accepted set narrower than the grid step. Regressions weighted by the
ACS person weight; standard errors clustered by household.
\end{minipage}
\end{table}

\paragraph{Same-age oldest pair.}
The first stage is large, 66 percentage points, because a same-age pair guarantees
a second child: the coefficient is approximately one minus the probability that a comparable
mother without such a pair already has a second coresident child. Mothers with a
same-age pair live in 0.33 more rooms (standard error 0.017) and 0.22 more bedrooms,
and they are 3.1 percentage points more likely to own. Scaling by the first stage,
a second child is associated with half a room, a third of a bedroom, and a
4.7-point higher ownership rate. The first stage falls from 52 points at event age
3 to 35 points at event age 5, as more mothers without a same-age pair have a second
child, while the room and bedroom reduced forms change little. The two-stage
estimates therefore rise with the event age, to 1.1 rooms and 15 points of ownership at age 5.
At event age 3 the ownership reduced form is 1.6 points and its confidence interval
includes zero. These magnitudes carry the qualifications above: the proxy mixes
twins with closely spaced and blended sibling pairs, and twin births are associated
with maternal characteristics that may affect housing directly.

\paragraph{Same-sex oldest pair.}
The first stage is 3.3 percentage points pooled over event ages 0 to 5 and rises to
4.7 points at event age 3 and 6.0 points at event age 5, consistent with more
third births as time passes since the second. With more than 650,000 mothers it is
precisely estimated, with an $F$ statistic of 715. The pooled figure is smaller
than the 6.2 points that \citet[p.~461]{angrist1998children} estimate in the 1980
Census, but the samples differ: ours ends five years after the second birth,
whereas theirs includes mothers whose second child is up to seventeen, and the
two censuses are a generation apart. Mothers whose two oldest
children share a sex live in 0.038 fewer rooms (standard error 0.005) and 0.034
fewer bedrooms (0.003), and they are 0.25 percentage points less likely to own, a
difference whose confidence interval includes zero. The room and bedroom
differences are larger at event ages 3 and 5, about 0.05 and 0.04 respectively, and
the ownership differences remain small and imprecise.

The two-stage estimate for rooms is $-1.15$, with an Anderson--Rubin set between
$-1.4$ and $-0.9$. It is large because the reduced form, though small, is divided by
a first stage of 0.033:
$-0.0376/0.0327=-1.15$. Taken literally, a third child would cost the family more
than a full room of living space. We do not interpret it this way. The estimate is
what one would obtain if same-sex siblings changed housing directly. Under the
exclusion restriction the reduced form equals $\pi\beta$, and any direct effect of
sibling sex on housing adds to it in full while the first stage stays the same, so
the direct effect enters the two-stage estimate multiplied by $1/\pi$, about 31.
\citet{jones2015exclusion} shows how such direct effects on mothers whose fertility
does not respond to the instrument bias the instrumental-variable estimand.
Bedrooms are the natural channel. Two children of the same sex can share a bedroom
more easily than a boy and a girl, so a family with same-sex children may need one
fewer bedroom at a given size, whether or not it has a third child.
\citet[p.~775]{angrist2010multiple} raise economies of scale in room or clothes
sharing as a potential violation of the same-sex exclusion restriction, attributing
the concern to \citet{rosenzweig2000natural}. They check it in Israeli subsamples
in which sex composition does not predict family size and find no reduced-form
relation between sex composition and their outcomes, which concern the children's
schooling and later lives rather than housing (p.~799). The pattern in Table~\ref{tab:acsiv_pooled}, in which the
bedroom difference accounts for nearly all of the room difference, is what this
channel predicts. It is consistent with the channel but does not establish it, and
we do not test the rooms and bedrooms differences against each other, which would
require their joint covariance.

\begin{table}[htbp]
\centering
\footnotesize
\caption{Estimates at event ages 3 and 5}
\label{tab:acsiv_ages}
\setlength{\tabcolsep}{4pt}
\begin{tabular}{@{}lrrrrrrrl@{}}
\toprule
 & & & Instrument & Reduced & First stage & First-stage & & AR accepted \\
Outcome & Mothers & Households & equal to one & form & (pp) & $F$ & 2SLS & grid values \\
\midrule
\multicolumn{9}{@{}l}{\textit{A. Same-age oldest pair, event age 3}} \\
Rooms & 147,986 & 147,922 & 2,726 & 0.332 & 52.46 & 49,032 & 0.633 & [0.5, 0.7] \\
 & & & & (0.042) & (0.24) & & (0.080) & \\
Bedrooms & 147,986 & 147,922 & 2,726 & 0.212 & 52.46 & 49,032 & 0.405 & \{0.4\} \\
 & & & & (0.023) & (0.24) & & (0.043) & \\
Owner (pp) & 147,986 & 147,922 & 2,726 & 1.63 & 52.46 & 49,032 & 3.11 & [0, 6] \\
 & & & & (1.14) & (0.24) & & (2.17) & \\
\addlinespace
\multicolumn{9}{@{}l}{\textit{A. Same-age oldest pair, event age 5}} \\
Rooms & 142,101 & 142,039 & 2,498 & 0.382 & 34.83 & 24,972 & 1.096 & [0.9, 1.3] \\
 & & & & (0.041) & (0.22) & & (0.119) & \\
Bedrooms & 142,101 & 142,039 & 2,498 & 0.234 & 34.83 & 24,972 & 0.671 & [0.6, 0.8] \\
 & & & & (0.023) & (0.22) & & (0.067) & \\
Owner (pp) & 142,101 & 142,039 & 2,498 & 5.26 & 34.83 & 24,972 & 15.11 & [10, 20] \\
 & & & & (1.15) & (0.22) & & (3.30) & \\
\addlinespace
\multicolumn{9}{@{}l}{\textit{B. Same-sex oldest pair, event age 3}} \\
Rooms & 108,847 & 108,825 & 54,561 & $-$0.051 & 4.67 & 174 & $-$1.089 & [$-$1.6, $-$0.6] \\
 & & & & (0.012) & (0.35) & & (0.284) & \\
Bedrooms & 108,847 & 108,825 & 54,561 & $-$0.040 & 4.67 & 174 & $-$0.856 & [$-$1.2, $-$0.6] \\
 & & & & (0.007) & (0.35) & & (0.165) & \\
Owner (pp) & 108,847 & 108,825 & 54,561 & $-$0.53 & 4.67 & 174 & $-$11.31 & [$-$26, 2] \\
 & & & & (0.36) & (0.35) & & (7.82) & \\
\addlinespace
\multicolumn{9}{@{}l}{\textit{B. Same-sex oldest pair, event age 5}} \\
Rooms & 92,246 & 92,229 & 46,201 & $-$0.053 & 6.04 & 214 & $-$0.874 & [$-$1.3, $-$0.5] \\
 & & & & (0.013) & (0.41) & & (0.235) & \\
Bedrooms & 92,247 & 92,230 & 46,201 & $-$0.045 & 6.04 & 214 & $-$0.745 & [$-$1.0, $-$0.5] \\
 & & & & (0.007) & (0.41) & & (0.137) & \\
Owner (pp) & 92,247 & 92,230 & 46,201 & $-$0.33 & 6.04 & 214 & $-$5.41 & [$-$18, 6] \\
 & & & & (0.40) & (0.41) & & (6.59) & \\
\bottomrule
\end{tabular}
\par\medskip
\begin{minipage}{0.97\linewidth}
\footnotesize\textit{Notes:} See Table~\ref{tab:acsiv_pooled}. Event age is the
oldest child's age in panel~A and the second-oldest child's age in panel~B, in
whole years at the interview.
\end{minipage}
\end{table}

\begin{figure}[htbp]
\centering
\includegraphics[width=\linewidth]{\acsivfigdir/acs_iv_reduced_forms.pdf}
\caption{Reduced-form differences in housing by sibling composition}
\label{fig:acsiv_rf}
\par\medskip
\begin{minipage}{0.97\linewidth}
\footnotesize\textit{Notes:} Each point is the reduced-form coefficient $\delta$ of
equation~\eqref{eq:acsiv_rf_fs}, with a 95 percent confidence interval from
household-clustered standard errors, for event ages 0--5 pooled and for event ages
3 and 5. Top row: mothers whose two oldest linked children are the same age,
relative to mothers whose oldest child is the same age but who have no same-age
pair. Bottom row: mothers whose two oldest linked children have the same sex,
relative to mothers of one boy and one girl. Ownership in percentage points.
\end{minipage}
\end{figure}

The two designs measure different objects, and their estimates should not be
averaged or set against each other. Under its assumptions, the same-age design says
that a second young child raises the space a family occupies and its likelihood of
owning. The same-sex design does not deliver a credible estimate for a third child,
because sibling sex appears to act on housing directly, through the composition of
the family and not only its size.
